\section{Conclusion}
\label{sec:conclusion}
%
We propose IDESplat, which iteratively applies warp operations and integrates multi-level epipolar attention maps to enhance depth probability estimation for accurate Gaussian mean prediction.
%
Specifically, we first design a depth probability boosting unit to amplify cross-view feature similarity in a multiplicative manner.
%
Then, we build an iterative depth estimation process by stacking multiple DPBUs, gradually identifying the most likely depth locations.
%
Additionally, during this iterative process, we progressively narrow the depth search range and increase feature size to achieve more precise depth estimates.
%
For the other Gaussian parameters, we design a gaussian focused module to select the most relevant Gaussian tokens for attention-based feature interaction.
%
Experimental results on large-scale benchmarks clearly show that IDESplat achieves excellent reconstruction quality and strong domain generalization ability.
%
% The superior performance of IDESplat can be attributed to its iteratively enhanced depth probability prediction design, which progressively refines the depth map.
%
Our current model still has limitations: although it enables real-time 3D reconstruction, its efficiency could be further improved, and it requires camera pose input as a prerequisite.
%
Potential directions for improvement include further enhancing inference speed and exploring a pose-free framework.



\section*{Acknowledgement}
% \vspace{2mm}
% This work was supported by the National Natural Science Foundation of China (No.~62476051).
{\raggedright
This work was supported by the National Natural Science Foundation of China under Grant No.~62476051.\par
}
